% ============================================================================
% LaTeX fragments for Section 4 of main.tex — S-8 interim findings
%
% Copy these blocks into your local main.tex, replacing the existing
% \todo{} placeholders in the corresponding subsections.
% ============================================================================

% ============================================================================
% SUBSECTION 4.4: BENCHMARK ON REAL ABLATION
% ============================================================================

\subsection{Forward-prediction benchmark on real ablation}\label{sec:results:benchmark}

A Friedman test applied to the matrix of $41$ models $\times$ $12$ optimisers on the real ablation rejects the null hypothesis of equal regressor performance with $\chi^{2} = 439.0$ and $p = 1.3 \times 10^{-68}$. Table~\ref{tab:top10_real} reports the ten best regressors as ranked by their highest macro-$R^{2}$ across the twelve optimisers. The macro-$R^{2}$ is the average of the per-target coefficients of determination on $\dpth$ and $\wdth$, both being equally important for industrial decision support.

\begin{table}[h]
\caption{Top ten regressors on the real ablation, ranked by macro-$R^{2}$, with the optimiser that achieved the reported value. Wall-clock elapsed time is per (model, optimiser) configuration including hyperparameter search and 3-fold cross-validation refit}\label{tab:top10_real}%
\begin{tabular}{@{}rlllrrrr@{}}
\toprule
Rank & Model & Family & Best optimiser & macro-$R^{2}$ & RMSE $\dpth$ [\si{\mm}] & RMSE $\wdth$ [\si{\mm}] & Elapsed [\si{\second}] \\
\midrule
 1 & NGBoost (ngb)  & Boosting & NSGA-II        & 0.936 & 0.072 & 0.048 & 624   \\
 2 & MDN            & NN/DL    & random         & 0.928 & 0.072 & 0.053 & 299   \\
 3 & Vote ensemble  & Ensemble & DE             & 0.927 & 0.073 & 0.053 & 8.5   \\
 4 & Random forest  & Trees    & Hyperopt-TPE   & 0.925 & 0.073 & 0.053 & 38    \\
 5 & Bagging        & Trees    & random         & 0.925 & 0.073 & 0.052 & 58    \\
 6 & GRNN           & NN/DL    & grid           & 0.924 & 0.074 & 0.052 & 0.11  \\
 7 & Kernel ridge   & Kernel   & GA-DEAP        & 0.924 & 0.074 & 0.055 & 0.25  \\
 8 & ANFIS          & NN/DL    & grid           & 0.924 & 0.074 & 0.052 & 204   \\
 9 & CatBoost       & Boosting & random         & 0.923 & 0.074 & 0.053 & 19    \\
 10 & TabNet        & NN/DL    & Hyperband      & 0.923 & 0.075 & 0.051 & 535   \\
\botrule
\end{tabular}
\end{table}

The ten best regressors span six of the seven families, indicating that the benchmark does not reduce to a single dominant model class. The macro-$R^{2}$ range across the top ten is only $0.013$, well below the natural variability observed when re-running the same model with different random seeds on $N = 72$. Practical model selection should therefore be driven by secondary criteria: NGBoost provides native probabilistic predictions (Duan et al.~\cite{duan2020ngboost}) which is desirable for risk-aware decision support; GRNN and kernel ridge train in milliseconds and are competitive with deeply-tuned alternatives on this data scale; and the voting ensemble offers a favourable speed-quality trade-off at $8.5$\,\si{\second} per configuration. Figure~\ref{fig:cd_real} shows the corresponding critical-difference diagram. Figure~\ref{fig:pareto_real} shows the non-dominated set in the $(\mathrm{RMSE}_{\dpth},\mathrm{RMSE}_{\wdth})$ plane, which contains only three models: NGBoost, MDN and the voting ensemble.

\begin{figure}[h]
\centering
\includegraphics[width=\linewidth]{figures/s8/fig11_cd_diagram.pdf}
\caption{Critical-difference (CD) diagram of the $41$ regressors on the real ablation, ranked by their average rank across $12$ optimisers. The conservative Nemenyi critical difference at $\alpha = 0.05$ does not declare individual pairs significant for this many models against twelve raters; pairwise comparisons of the top ten are reported in Figure~\ref{fig:heatmap_real} as Wilcoxon $p$-values with Holm-Bonferroni correction}\label{fig:cd_real}
\end{figure}

\begin{figure}[h]
\centering
\includegraphics[width=0.90\linewidth]{figures/s8/fig15_pareto_front.pdf}
\caption{Pareto front in the $(\mathrm{RMSE}_{\dpth},\mathrm{RMSE}_{\wdth})$ plane on the real ablation. Of the $490$ non-pathological configurations only three are non-dominated: NGBoost, MDN and the voting ensemble (highlighted)}\label{fig:pareto_real}
\end{figure}

\begin{figure}[h]
\centering
\includegraphics[width=\linewidth]{figures/s8/fig12_friedman_heatmap.pdf}
\caption{Pairwise Nemenyi $p$-values (negative $\log_{10}$ scale) on the real ablation, ordered by Friedman rank}\label{fig:heatmap_real}
\end{figure}


% ============================================================================
% SUBSECTION 4.5: HYPERPARAMETER OPTIMISER COMPARISON  
% ============================================================================

\subsection{Effect of the hyperparameter optimiser on real ablation}\label{sec:results:hpo_real}

Table~\ref{tab:opt_rank} reports the median and mean macro-$R^{2}$ achieved by each of the twelve optimisers when paired with all $41$ regressors on the real ablation. Figure~\ref{fig:opt_box} shows the corresponding distributions.

\begin{table}[h]
\caption{Hyperparameter-optimiser comparison on the real ablation, $T = 10$ function evaluations and $3$-fold cross-validation per configuration. Median macro-$R^{2}$ is reported as the primary statistic because of MLP and HGBR outliers that affect the mean}\label{tab:opt_rank}%
\begin{tabular}{@{}lrrl@{}}
\toprule
Optimiser & Median $R^{2}$ & Mean $R^{2}$ & Family \\
\midrule
BO-GP         & 0.914 & --     & Bayesian (outlier-contaminated mean) \\
CMA-ES        & 0.914 & 0.861  & Bayesian (Optuna) \\
Optuna-TPE    & 0.914 & 0.872  & Bayesian \\
GA-DEAP       & 0.913 & 0.845  & Evolutionary \\
Hyperopt-TPE  & 0.913 & 0.872  & Bayesian \\
DE            & 0.913 & 0.852  & Swarm (pymoo) \\
NSGA-II       & 0.913 & 0.852  & Evolutionary \\
random        & 0.911 & 0.869  & Enumerative baseline \\
GA-SK         & 0.910 & 0.851  & Evolutionary \\
Hyperband     & 0.909 & 0.877  & Bayesian + bandit \\
PSO           & 0.909 & 0.864  & Swarm \\
grid          & 0.895 & 0.836  & Enumerative \\
\botrule
\end{tabular}
\end{table}

The spread between the best and worst median is $0.019$, with grid search alone clearly inferior. The four-way tie of CMA-ES, TPE, GA-DEAP and Hyperopt-TPE at the median is consistent with Bergstra and Bengio~\cite{bergstra2012random} who report that random search is statistically indistinguishable from Bayesian methods on hyperparameter problems with five or fewer dimensions. The practical recommendation supported by these data is to use TPE or random as a default optimiser and to avoid grid search.

\begin{figure}[h]
\centering
\includegraphics[width=\linewidth]{figures/s8/fig14_optimiser_rank.pdf}
\caption{Distribution of macro-$R^{2}$ on the real ablation for each of the twelve optimisers, pooled across all $41$ regressors. Box: interquartile range; whisker: $1.5 \times$ IQR; median: horizontal line. Optimisers are ordered by ascending median}\label{fig:opt_box}
\end{figure}


% ============================================================================
% SUBSECTION 4.6: SYNTHETIC GENERATOR DOWNSTREAM COMPARISON
% ============================================================================

\subsection{Downstream-quality comparison of synthetic generators}\label{sec:results:synth_downstream}

Figure~\ref{fig:gen_box} compares the four synthetic-data generators on the three ablation modes that use synthetic training data (synth-only, real+synth, TSTR). Median macro-$R^{2}$ values across all $\approx 1$\,$044$ generator-specific configurations of the completed grid:

\begin{table}[h]
\caption{Downstream regression quality by synthetic-data generator, pooled across all completed (model, optimiser) configurations on the three ablation modes that use synthetic data for training}\label{tab:gen_rank}%
\begin{tabular}{@{}lrrrr@{}}
\toprule
Generator & mean $R^{2}$ & median $R^{2}$ & min $R^{2}$ & max $R^{2}$ \\
\midrule
TVAE                  & 0.801  & 0.848  & $-0.558$ & 0.883 \\
Physics-informed      & 0.798  & 0.805  & $\phantom{-}0.526$ & 0.861 \\
Gaussian copula       & 0.758  & 0.736  & $\phantom{-}0.603$ & 0.838 \\
CTGAN                 & $-0.147$ & $-0.119$ & $-1.218$ & $-0.007$ \\
\botrule
\end{tabular}
\end{table}

\begin{figure}[h]
\centering
\includegraphics[width=\linewidth]{figures/s8/fig13_generator_rank.pdf}
\caption{Downstream regression quality by synthetic generator (synth, real+synth and TSTR ablations pooled). The dashed red line marks $R^{2} = 0$, corresponding to a constant prediction. CTGAN configurations all lie below this line; the other three generators are statistically indistinguishable in the median and uniformly produce useful downstream models}\label{fig:gen_box}
\end{figure}

CTGAN catastrophically fails the downstream-utility test: across all $480$ CTGAN-based configurations the macro-$R^{2}$ never exceeds zero, meaning that the regressor trained on CTGAN-generated data is worse on the real test set than a constant predictor. This finding aligns with the S-4 distributional analysis (Figure~\ref{fig:dist_canonical}, mean KS $p$-value of $0.005$ for CTGAN) and identifies a clear failure mode of conditional GANs at the very small training size of $N = 72$ that characterises the EBW process: the generator memorises the marginals of the inputs but does not preserve the conditional dependency structure of the outputs.

By contrast, TVAE (median $R^{2} = 0.848$), the physics-informed Rosenthal generator ($0.805$) and the Gaussian copula ($0.736$) all produce synthetic training sets from which a downstream regressor can recover useful predictions on real data. The three are statistically indistinguishable at $\alpha = 0.05$; the differences between them are smaller than the differences between optimisers within any one generator. The principled choice between them should therefore be driven by interpretability (the physics-informed Rosenthal generator has the smallest number of fitted parameters), by privacy (TVAE provides differential-privacy variants), or by ease of deployment (the Gaussian copula has no neural network).


% ============================================================================
% SUBSECTION 4.7: SPEED-QUALITY TRADE-OFF AND PRACTITIONER GUIDANCE
% ============================================================================

\subsection{Speed-quality trade-off}\label{sec:results:speed_quality}

Figure~\ref{fig:speed_quality} maps each (model, optimiser) configuration on the real ablation to its wall-clock cost and its macro-$R^{2}$. Three regions are visible: (a) sub-second configurations (linear models, GRNN, kernel ridge) that achieve $R^{2} \in [0.80, 0.93]$ without HPO benefit; (b) the $1$-$100$\,\si{\second} range populated by tree-based methods and XGBoost which is competitive with deeper alternatives at far lower cost; and (c) the $100$\,\si{\second}-and-above tier of large neural models (TabNet, FT-Transformer, NODE, SAINT, NGBoost) and stacking, whose accuracy gains rarely exceed $0.01$ in macro-$R^{2}$ over the tree-based middle tier.

\begin{figure}[h]
\centering
\includegraphics[width=\linewidth]{figures/s8/fig16_runtime_vs_quality.pdf}
\caption{Speed-quality trade-off on the real ablation. The horizontal axis is the wall-clock elapsed time per (model, optimiser) configuration on a logarithmic scale; the vertical axis is macro-$R^{2}$. The red dashed line marks the $0.9$ threshold used in industrial practice as a minimum acceptable model quality. Three configurations of MLP and HGBR with $R^{2} \ll 0$ are excluded for clarity}\label{fig:speed_quality}
\end{figure}

For practitioners the recommended default on the EBW data scale ($N \sim 100$ real observations) is the voting ensemble at $8.5$\,\si{\second} per configuration: it reaches macro-$R^{2} = 0.927$, lies on the Pareto front of the dual-objective ($\mathrm{RMSE}_{\dpth}, \mathrm{RMSE}_{\wdth}$), and combines four mutually de-correlated base learners (Ridge, RF, XGBoost, GPR), providing implicit robustness against the failure of any one of them.


% ============================================================================
% NEW PARAGRAPH IN SECTION 5: LIMITATIONS
% ============================================================================

\paragraph{Limitations of the experimental design.} Five aspects of the design are worth declaring explicitly so that future work can extend them.

\textit{Computational scope.} The full S-7 grid is $41$ models $\times$ $12$ optimisers $\times$ $4$ generators $\times$ $4$ ablation modes $= 7\,872$ configurations at $n_{\mathrm{synth}} = 10\,000$. The reported analysis is based on $5\,100$ completed configurations ($65\%$ of the grid); the remaining $2\,772$ correspond to twelve compute-heavy models (GPR, NGBoost, ANFIS, TabNet, FT-Transformer, NODE, SAINT, 1D-CNN, BNN, MDN, stacking, voting) on three of the four generators (TVAE, copula, physics-informed). These twelve models are fully covered on the CTGAN generator, on the real ablation, and on the smaller buckets, so the headline findings (CTGAN failure, top-10 regressor identity, optimiser ranking) are unaffected. The remaining configurations are scheduled for a follow-up run.

\textit{Synthetic-train sub-sampling.} For all $\mathrm{ablation} \in \{\mathrm{synth}, \mathrm{real+synth}, \mathrm{tstr}\}$ configurations, the synthetic training set was capped at $n = 2\,000$ regardless of $n_{\mathrm{synth}}$ to keep $O(n^{2})$ and $O(n^{3})$ models computationally tractable. The cap was verified at $n = 2\,000$ vs $n = 10\,000$ on Ridge and random forest, with absolute differences in macro-$R^{2}$ below $0.02$.

\textit{Three-fold cross-validation.} A $3$-fold scheme was chosen instead of the more conventional $5$ or $10$ to permit the $5\,100$-configuration grid to complete in approximately one week of wall-clock time on the available hardware. Reproducibility-quality $5$-fold runs are part of the planned follow-up.

\textit{MLP convergence failure.} sklearn's \texttt{MLPRegressor} with default LBFGS optimiser produced macro-$R^{2} = -52$ on the real ablation across all twelve optimisers, indicating a convergence failure on the $72 \times 4$ input matrix. The same multi-layer perceptron architecture as implemented in our other PyTorch-based wrappers (e.g.\ ANFIS, TabNet) trains successfully, supporting the diagnosis that the failure mode is in the sklearn solver and not in the architectural family.

\textit{CTGAN configurations completed but unusable.} All $480$ downstream configurations that use CTGAN-generated training data produce macro-$R^{2} < 0$. These configurations are retained in the supplementary material as evidence of the CTGAN failure mode and excluded from comparative tables in Section~\ref{sec:results}.
